\section{Related Literature and Contribution Boundary}

The paper sits at the intersection of three literatures: compatibility and endogenous product design, innovation incentives under standards and compatibility, and the allocation of innovative effort across technological activities.

First, compatibility need not be welfare improving once product design is endogenous. \citet{Boom2001} studies a three-stage compatibility--product-specification--pricing game and shows that endogenous product selection can overturn the welfare gains from compatibility. \citet{Ruiz2004} studies international standardization policy with differentiated products and shows that endogenous product characteristics affect the consequences of recognizing foreign standards. These papers establish that firms can redesign products in response to compatibility policy. Our adjustment margin is different: product substitutability is fixed, while firms reallocate a scarce R\&D budget between a common technological layer and a proprietary layer.

Second, a substantial literature links compatibility to innovation incentives. \citet{BondSmith2019} studies compatibility, endogenous R\&D, and commercialization in markets with network effects. \citet{MaruyamaZennyo2017} analyzes process innovation and endogenous application compatibility, showing that improved R\&D efficiency can induce lower compatibility and reduce welfare. \citet{HeywoodWangYe2022} compares R\&D rivalry and cooperation when compatibility is endogenous. \citet{Llanes2024} studies incentives to develop technologies for inclusion in a technical standard and shows that standardization can generate either insufficient or excessive innovation. In a dynamic general-equilibrium setting, \citet{AcemogluGanciaZilibotti2012} analyzes the competing roles of innovation and standardization and characterizes an optimal speed of standardization. These contributions make clear that neither ``compatibility affects innovation'' nor ``partial standardization may be optimal'' is novel by itself.

Third, the allocation of innovative effort across alternative activities is also established. \citet{KretschmerReitzig2013} model firms' R\&D diversification in emerging standard settings, where firms trade off cooperation on a joint standard against own activities. \citet{daSilva2024} studies scarce managerial resources allocated across internal and external R\&D and shows how spillovers and substitutability across research approaches affect the allocation. Evidence on modular product innovation likewise suggests that external interface standards can redirect innovative attention toward internal interfaces and modules \citep{ChenLiu2005}.

The contribution of the present paper is therefore deliberately narrower than any of those broad statements. We study a regulator that chooses a continuous standardization scope before firms allocate a fixed innovation capacity. The same scope parameter enters private and social R\&D returns with opposite strategic signs: greater transferability lowers the private relative return to common innovation because it strengthens rivals, while increasing its social relative return because the innovation serves more products. This produces a full-game policy reversal. Holding R\&D composition fixed, complete standardization is optimal; complete standardization is also first-best. Yet once firms endogenously redirect a fixed R\&D portfolio, sufficiently strong product-market rivalry makes the regulator choose a unique interior scope. The novel object is this fixed-allocation/full-standardization to endogenous-portfolio/selective-standardization comparison and its threshold characterization, not the component mechanisms considered separately.
